\chapter{Gas path analysis of the CFM56-7B26, in practice}
\label{ch:gpa-practice}

\begin{quote}\itshape
This chapter in brief: \cref{ch:background} gave the theory of gas path analysis and
\cref{ch:model} built the influence matrix. Here the two meet the engine: known faults are
injected into the nonlinear twin, the measurements are inverted exactly as a fielded GPA
system would invert them, and we measure what comes back --- how much of the fault, spread
over which components, how small a fault can still be named, and what happens when a real
engine carries eight of them at once. This is the classical GPA performance study for this
specific engine, and it is the quantitative floor every later chapter is measured against.
\end{quote}

\keyidea{Run the textbook GPA on this engine and the cockpit sensor set returns, on average, a
quarter of the injected fault with $85\,\%$ of the estimated magnitude smeared onto healthy
components --- the wall is visible before any learning, any fleet, or any hypothesis. Station
probes turn most of that around, which is why the sensor conclusion recurs in every later
chapter.}


\section{The experiment: honest forward, classical inverse}
\label{sec:gpa-setup}

A GPA study is only as honest as its forward model. Inverting the same linear matrix that
generated the measurements would be circular --- the answer would be exact by construction and
would say nothing about the real thing. So the two directions are deliberately different
instruments:

\begin{description}
  \item[Forward (truth).] The \textbf{nonlinear} neural twin of \cref{sec:f8-l1}: a health
    state goes in, the eight gas-path channels come out with the full curvature of the pyCycle
    physics, $13$--$19\times$ closer to it than a linearization. This plays the engine.
  \item[Inverse (the diagnosis).] The \textbf{linear} regularized weighted least squares of
    \cref{eq:wls} on the F1 cruise influence matrix, at the engineering defaults registered in
    \cref{ch:decisions} (prior $\sigma = 2\,\%$, $\lambda = 1$). This plays the monitoring
    system --- and, like every fielded GPA, it does not know the physics is curved.
  \item[Sensors.] Gaussian noise per channel at the magnitudes of the versioned fault catalog,
    expressed as a percentage of the healthy cruise value: N2 $0.05\,\%$, $W_f$ $0.50\,\%$,
    EGT $0.22\,\%$, P25 $0.40\,\%$, T25 $0.42\,\%$, PS3 $0.40\,\%$, T3 $0.28\,\%$. Every
    experiment is run twice, once with the cockpit set (N2, $W_f$, EGT) and once with the
    extended set that adds the four station probes.
\end{description}

All of it is \code{scripts/gpa_study.py}; every number below is read from its artifact
(\code{data/processed/gpa/gpa_study.json}), never transcribed.

\paragraph{One run, end to end.} Take the canonical fault: the HPC loses $1\,\%$ of
efficiency, nothing else. The twin says the cockpit sees N2 fall by $0.12\,\%$, fuel flow rise
by $0.57\,\%$ and EGT rise by $0.52\,\%$. Feed that to the inversion and it returns
$-0.26\,\%$ on HPC efficiency --- a quarter of the truth --- with the rest of the magnitude
scattered across the other nine parameters. Repeat the run 500 times with fresh sensor noise
and the diagnosis names the HPC correctly in $22\,\%$ of them; its favourite wrong answer is
the HPT. Nothing has gone wrong: this is what the textbook method does on this engine with
these sensors, and the rest of the chapter is the systematic version of that one run.

\section{The fault signature atlas}
\label{sec:gpa-atlas}

Before any inversion, the forward half deserves to be seen. \Cref{fig:gpa-atlas} is the
classical GPA deliverable --- one panel per health parameter, showing how each measured channel
responds to a $+1\,\%$ deviation of it. It is the same information as the heat map of
\cref{fig:icm}, drawn the way a practitioner reads it: as ten fingerprints to be told apart.

\begin{figure}[htbp]
  \centering
  \includegraphics[width=\textwidth]{gpa_atlas}
  \caption{Fault signature atlas at cruise: the response of every measured channel to a
  $+1\,\%$ deviation of each health parameter. Blue bars are the cockpit channels, grey the
  station probes an operator may not have. Read across the blue bars only and several panels
  become hard to tell apart --- that is the diagnosis problem in one picture.
  Auto-generated.}
  \label{fig:gpa-atlas}
\end{figure}

Two readings matter. First, \emph{shape}: covering the grey bars, the HPC-efficiency and
HPT-efficiency panels are near-copies of one another (N2 $+0.12$ vs $+0.16$, $W_f$ $-0.57$ vs
$-0.70$, EGT $-0.52$ vs $-0.68$ \%\,per\,\%) --- proportional to within a couple of degrees of
angle, the pair \cref{sec:icm} measured at $1.3^{\circ}$. Second, \emph{size}: the panels are
not drawn to the same effective scale. Fan flow capacity moves the cockpit channels seventeen
times harder than HPT flow capacity does; a turbine whose effective nozzle area has opened by
$1\,\%$ barely disturbs the three numbers an airline receives. The station probes are what
break both problems --- PS3 alone responds to HPT flow capacity at $-1.05\,\%$ per \%, an
order of magnitude more than any cockpit channel.

\section{What the inversion gives back}
\label{sec:gpa-recovery}

Inject each fault at $1\,\%$ in its physically expected direction, invert, and compare
(\cref{fig:gpa-recovery} and \cref{tab:gpa-study}). Two quantities matter to a maintenance
engineer: how much of the magnitude comes back, and how much of what comes back is attached to
the wrong component --- the \emph{smearing index}, the fraction of $\|\hat{\mathbf{x}}\|_1$
sitting on healthy parameters.

\begin{figure}[htbp]
  \centering
  \includegraphics[width=\textwidth]{gpa_recovery}
  \caption{Single-fault recovery at $1\,\%$, cockpit versus extended sensors. Left: the
  fraction of the injected magnitude the estimator returns on the correct parameter (noise-free
  --- so this is the ceiling, not an average case). Right: the smearing index, the share of the
  estimated magnitude landing on components that are in fact healthy. Auto-generated.}
  \label{fig:gpa-recovery}
\end{figure}

\begin{table}[htbp]
  \centering\small
  \caption{The GPA scorecard for this engine, per health parameter and sensor set: magnitude
  recovered, smearing index, correct-isolation rate at a $1\,\%$ fault, and the smallest fault
  the nearest-signature rule names correctly in $90\,\%$ of noisy draws. Auto-generated from
  \texttt{scripts/gpa\_study.py}.}
  \label{tab:gpa-study}
  % AUTOGENERATED by scripts/make_report_assets.py — do NOT edit by hand
\begin{tabular}{@{}lrrrrrrrr@{}}
\toprule
& \multicolumn{2}{c}{recovered [\%]} & \multicolumn{2}{c}{smearing} & \multicolumn{2}{c}{isolation} & \multicolumn{2}{c}{min.\ detectable [\%]} \\
\cmidrule(lr){2-3}\cmidrule(lr){4-5}\cmidrule(lr){6-7}\cmidrule(lr){8-9}
Parameter & ckpt & ext & ckpt & ext & ckpt & ext & ckpt & ext \\
\midrule
$\eta_{fan}$ & 32 & 41 & 0.84 & 0.87 & 0.31 & 0.52 & $>3$ & 3 \\
$\Gamma_{fan}$ & 76 & 79 & 0.68 & 0.80 & 0.67 & 0.99 & 2 & 0.75 \\
$\eta_{LPC}$ & 15 & 39 & 0.88 & 0.84 & 0.17 & 0.22 & $>3$ & $>3$ \\
$\Gamma_{LPC}$ & 27 & 69 & 0.84 & 0.82 & 0.31 & 0.90 & $>3$ & 1 \\
$\eta_{HPC}$ & 26 & 89 & 0.86 & 0.76 & 0.22 & 0.73 & $>3$ & 1.5 \\
$\Gamma_{HPC}$ & 46 & 79 & 0.79 & 0.77 & 0.16 & 0.10 & $>3$ & $>3$ \\
$\eta_{HPT}$ & 45 & 70 & 0.80 & 0.81 & 0.09 & 0.81 & $>3$ & 1.5 \\
$\Gamma_{HPT}$ & 2 & 90 & 0.98 & 0.75 & 0.06 & 0.79 & $>3$ & 1.5 \\
$\eta_{LPT}$ & 21 & 32 & 0.87 & 0.91 & 0.22 & 0.58 & $>3$ & $>3$ \\
$\Gamma_{LPT}$ & 15 & 32 & 0.91 & 0.91 & 0.44 & 0.81 & $>3$ & 1.5 \\
\bottomrule
\end{tabular}

\end{table}

With the cockpit set the recovered fraction spans $2$--$76\,\%$ and the median smearing index
is $0.85$: on a typical single fault, six sevenths of the diagnosed magnitude is attributed to
components that are fine. The extended set lifts recovery to $32$--$90\,\%$ while barely
moving the smearing index ($0.82$) --- more sensors buy \emph{magnitude}, not tidiness, because
the estimator still spreads what it cannot separate.

The extremes are the instructive part. \textbf{HPT flow capacity} is the pathological case:
$1.9\,\%$ of it comes back from the cockpit --- the fault is, for practical purposes, invisible
--- and $90\,\%$ from the extended set, a $47\times$ improvement bought by one probe at the HPC
exit. \textbf{HPC efficiency}, the parameter an operator most wants, goes from $26\,\%$ to
$89\,\%$. At the other end, the two LPT parameters improve barely at all ($21\to32\,\%$ and
$15\to32\,\%$): no probe in either set sits downstream of the low-pressure turbine, so no
sensor package considered here can see it. The certificate of \cref{ch:breakthrough} states
the same fact as a Fisher-information calculation; here it is simply measured.

\paragraph{Why nothing reaches $100\,\%$, even without noise.} The regularization is not a
defect but a choice, and it costs magnitude. The $\lambda\mathbf{P}_0^{-1}$ term of
\cref{eq:wls} pulls the estimate toward zero wherever the measurements do not insist
otherwise, so an unobservable direction returns almost nothing and an observable one returns
most of the truth. That shrinkage is exactly what keeps the ill-conditioned inversion from
exploding into the noise-amplified nonsense of the toy example in \cref{sec:gpa}; the price is
paid in bias.

\paragraph{A bias the certificate had to learn about.} The trade is visible in the numbers.
Under cockpit sensors the estimator's \emph{spread} is small ($0.02$--$0.35\,\%$) while its
\emph{bias} is large ($0.27$--$0.98\,\%$); total error stays between $0.44$ and $0.98\,\%$,
comfortably \emph{below} the Cram\'er--Rao bound of $1.04$--$1.98\,\%$. There is no paradox: the
Cram\'er--Rao bound constrains \emph{unbiased} estimators, and this one is deliberately biased,
so it is allowed to beat the bound in mean-square error by giving up accuracy of the mean.
This is the same discovery that refuted hypothesis H10.2 in \cref{sec:bt-honesty} --- the bound
gives the right \emph{shape} of the uncertainty and the wrong \emph{scale} --- reached here from
the opposite direction, on a single snapshot rather than an accumulated history.

\section{How small a fault can be named}
\label{sec:gpa-detectability}

Recovery measures magnitude; an operator also wants to know when a diagnosis can be
\emph{trusted}. Sweeping the injected magnitude and asking the classical nearest-signature rule
(\cref{sec:trad-overview}) to name the fault gives the detectability curves of
\cref{fig:gpa-detectability}: the smallest fault named correctly in $90\,\%$ of noisy draws is
the practical resolution limit of the sensor set.

\begin{figure}[htbp]
  \centering
  \includegraphics[width=\textwidth]{gpa_detectability}
  \caption{Correct-isolation rate against fault magnitude for four representative parameters,
  cockpit versus extended sensors; the dashed line is the $90\,\%$ criterion. The extended set
  crosses it between $0.75$ and $1.5\,\%$ for most faults; the cockpit set mostly never does
  within the $3\,\%$ range swept. Auto-generated.}
  \label{fig:gpa-detectability}
\end{figure}

The result is stark. Across the ten parameters the cockpit set reaches the $90\,\%$ criterion
for \textbf{one} of them (fan flow capacity, at $2\,\%$); the extended set reaches it for
\textbf{seven}, several as low as $0.75$--$1.5\,\%$. Three parameters --- LPC efficiency, HPC
flow capacity and LPT efficiency --- are never resolved by either set inside the swept range;
for those the classical rule usually abstains rather than guessing wrong, which is the correct
behaviour and worth more than a confident error. Set against the fleet's own fault magnitudes
($0.3$--$1.1\,\%$ after the v1.1 difficulty gate, \cref{sec:fleet-v11}), the reading is that
\emph{most faults in this benchmark sit below the cockpit set's resolution limit} --- which is
precisely why the isolation contest of \cref{sec:res-h2} ends in a tie near chance level, and
why it was pre-registered as a hypothesis about information rather than about algorithms.

\section{A real engine, eight faults at once}
\label{sec:gpa-multifault}

Single-fault studies are the textbook exercise; real engines deteriorate everywhere at once.
\Cref{fig:gpa-multifault} takes engine 99 of the fleet at cycle $19\,489$ --- late life, still
on wing --- and asks both sensor sets to diagnose its true ten-component state.

\begin{figure}[htbp]
  \centering
  \includegraphics[width=\textwidth]{gpa_multifault}
  \caption{A real fleet engine's end-of-life health state (black) and the GPA diagnosis of it
  under each sensor set, with the sensor-noise spread as error bars. The cockpit inversion
  recovers the hot-section efficiency loss but almost entirely misses the turbine flow-capacity
  opening that accompanies it; the extended set recovers both. Auto-generated.}
  \label{fig:gpa-multifault}
\end{figure}

The engine's true state is dominated by hot-section deterioration: HPT efficiency
$-4.20\,\%$ with the flow capacity opened $+2.39\,\%$, plus HPC efficiency $-2.33\,\%$ from
fouling and wear. The cockpit diagnosis gets the headline right --- HPT efficiency
$-3.28\,\%$, HPC efficiency $-2.47\,\%$ --- and then fails on the physics that would confirm it:
the turbine flow opening comes back as $+0.61\,\%$ of the true $+2.39$, so the classic
eroded-nozzle fingerprint is not visible in the diagnosis at all. The extended set recovers it
($+2.49\,\%$) and cuts the overall error from $1.05$ to $0.85\,\%$ RMSE.

Both sets share the same failure, and it is the honest lesson of the chapter: they invent
deterioration that is not there. True fan efficiency loss is $-0.46\,\%$; both diagnoses report
three times that. True LPT flow capacity is $+0.14\,\%$; both report about $-1.3\,\%$, wrong in
sign as well as magnitude. The estimator must place the measured residual \emph{somewhere}, and
in directions the sensors cannot resolve it places it wherever the regularization finds
cheapest. A maintenance decision taken on those two numbers would open a healthy module.

\section{What this chapter fixes in place}

Everything that follows is measured against this baseline, and three of its conclusions recur
so often that it is worth naming them here:

\begin{enumerate}
  \item \textbf{The wall is in the sensors, and it is visible before any learning.} With the
    cockpit set the textbook method recovers a quarter of a typical fault, smears $85\,\%$ of
    what it finds onto healthy components, and resolves one parameter in ten. No estimator
    choice in this document changes that, because \cref{sec:f8-lh2} shows the missing quantity
    is information, not algorithm.
  \item \textbf{Station probes are worth more than they cost.} Recovery, isolation and
    resolution all improve by roughly a factor of three, and the parameters they rescue
    (HPC efficiency, HPT flow capacity) are exactly the ones a shop visit is scheduled around.
    \Cref{sec:bt-honesty} turns that into a per-sensor purchasing calculation.
  \item \textbf{A GPA number without a confidence statement is not a diagnosis.} The same
    inversion returns a reliable HPC efficiency under one sensor set and reads pure noise under
    another, with nothing in its output to distinguish the two cases. That is the gap the
    certificate of \cref{ch:breakthrough} exists to close.
\end{enumerate}
